\section{从一个 336 字节请求对象开始：分配器究竟要交付什么}

前文已经说明页级 buddy 与对象级 slab 的基本分工。现在用一项具体任务贯穿慢路径和释放
路径：块层要创建请求 R37。教学化的 R37 占 336 字节，要求 64 字节对齐，保存操作类型、
逻辑块范围、缓冲区引用、完成状态和等待队列。提交线程可以睡眠，但设备完成中断不能睡眠；
请求发布到设备以后，局部函数即使返回，R37 仍必须存在。

朴素的 \code{allocate(336)} 接口没有表达以下问题：

\begin{enumerate}
\tightlist
\item 返回的是虚拟连续、物理连续还是仅由页表拼成的地址？
\item 分配失败时能否睡眠、回收页、启动 I/O 或进入文件系统？
\item 对象属于哪个 NUMA 节点和资源组，能否使用系统应急储备？
\item 内存清零、构造函数、对齐和调试红区由谁完成？
\item 调用者在何时可以发布指针，最后一个引用消失后何时才能复用槽位？
\end{enumerate}

因此，分配结果不只是一个地址。它同时兑现“容量、布局、上下文与生命周期”四类契约。

\Needspace{16\baselineskip}
\begin{longtable}{|L{0.18\linewidth}|L{0.27\linewidth}|L{0.31\linewidth}|L{0.14\linewidth}|}
\hline
契约 & R37 的要求 & 分配器或调用者保存的证据 & 失败结果 \\ \hline
\endfirsthead
\hline
契约 & R37 的要求 & 分配器或调用者保存的证据 & 失败结果 \\ \hline
\endhead
容量与对齐 & usable size 至少 336 B，按 64 B 对齐 &
cache object size、alignment、slab layout & 返回空指针或错误 \\\tablerowrule
地址性质 & CPU 可访问；请求对象本身不要求物理连续 &
虚拟地址、所属 slab、承载页 & 改用合适接口 \\\tablerowrule
执行上下文 & 提交路径可等待；中断完成路径不可进入直接回收 &
allocation context、GFP 约束、预留池身份 & 降级或消费预留 \\\tablerowrule
资源归属 & 记入发起任务所在 memcg 与 NUMA 策略 &
charge owner、node、cache 统计 & 局部回收或 memcg OOM \\\tablerowrule
生命周期 & 发布后保持到完成、引用归零和延迟读者结束 &
refcount、state、RCU callback、debug generation & 延迟释放或报告错误 \\
\hline
\end{longtable}

本章后续把这五类证据逐层建立。只有在最后一节，R37 才从“需要 336 字节”变成一个可以
安全提交、完成和回收的内核对象。

\section{可分配页从哪里来：固件范围不能直接当作空闲链表}

机器启动时，固件或平台描述提供若干物理地址范围。范围可能包含可用 RAM、固件保留区、
设备窗口、内核映像、启动页表和早期分配结果。此时还没有完整 buddy；若直接把所有
“RAM 范围”挂入空闲链表，分配器可能把正在保存内核代码或页表的页面交给普通请求。

早期分配器先维护“可用 extent 减去保留 extent”的区间集合。建立每页描述符、zone 和
buddy 元数据以后，系统再逐段释放确实可分配的页。这个转换具有提交意义：某个 PFN 只有在
固件类型、内核占用、体系结构保留和早期分配四项检查都通过后，才进入普通页分配器。

\begin{longtable}{|L{0.2\linewidth}|L{0.31\linewidth}|L{0.39\linewidth}|}
\hline
启动状态 & 代表字段 & 此阶段允许的动作 \\ \hline
固件内存图 & start、length、type、attributes & 识别候选 RAM 和设备/保留窗口 \\\tablerowrule
早期 extent 集 & available、reserved、allocation cursor &
在 buddy 建立前分配连续启动结构 \\\tablerowrule
页描述符数组 & PFN、flags、refcount、node、zone & 为每个受管页建立稳定身份 \\\tablerowrule
zone 初始化状态 & start PFN、spanned、present、managed pages、watermarks &
定义区域边界和可参与分配的容量 \\\tablerowrule
buddy 空闲表 & order、migratetype、free block head、count &
接受普通页分配与释放 \\
\hline
\end{longtable}

\begin{pdfdiagram}[scale=0.9]
  \node[os state, minimum width=34mm] (firmware) at (-56mm,18mm)
    {固件范围\\可用 RAM 与保留窗口};
  \node[os memory, minimum width=34mm] (early) at (-19mm,18mm)
    {早期 extent\\扣除内核、页表和启动分配};
  \node[os compute, minimum width=34mm] (pages) at (19mm,18mm)
    {页描述符\\PFN、zone、refcount};
  \node[os control, minimum width=34mm] (buddy) at (56mm,18mm)
    {zone/buddy\\按 order 与迁移类型入表};
  \draw[os bus] (firmware.east) -- (early.west);
  \draw[os bus] (early.east) -- (pages.west);
  \draw[os strong bus] (pages.east) -- (buddy.west);

  \node[os label, text width=34mm] at (-19mm,-4mm)
    {保留范围不能被普通请求取得};
  \node[os label, text width=36mm] at (39mm,-4mm)
    {只有 managed pages 进入分配容量};
  \draw[os guide] (-19mm,7mm) -- (-19mm,-12mm);
  \draw[os guide] (39mm,7mm) -- (39mm,-12mm);
\end{pdfdiagram}
\diagramnote{“物理地址属于 RAM”只说明介质类型，不说明页面当前空闲。启动转换要先扣除所有占用，再把具有页描述符身份的 managed pages 交给 buddy。}

\Needspace{8\baselineskip}
\subsection{页描述符为何独立于页内字节}

物理页保存实际数据，页描述符保存内核管理事实。代表性字段包括：

\begin{longtable}{|L{0.2\linewidth}|L{0.28\linewidth}|L{0.42\linewidth}|}
\hline
字段类别 & 代表字段 & 作用与边界 \\ \hline
身份与位置 & PFN、zone、node & 说明页面位于何处，不表示当前用途 \\\tablerowrule
生命状态 & flags、refcount、mapcount & 判断是否仍有引用、映射或特殊状态 \\\tablerowrule
用途关联 & mapping、index、private & 把页关联到文件偏移、匿名映射或子系统私有状态 \\\tablerowrule
分配器状态 & buddy order、migratetype、slab marker &
字段含义随页面当前状态变化，不能同时解释为多种用途 \\\tablerowrule
链表关系 & lru、buddy list、slab list & 同一嵌入节点在不同状态下由不同所有者使用 \\
\hline
\end{longtable}

稳定语义是“每个受管页有可定位的管理身份”；具体字段复用和布局属于实现选择。页面处于
buddy 空闲状态时，描述符中的 order 才能解释为连续空闲块阶数；页面交给 slab 后，相同存储
位置可能改作 slab 元数据。状态转换必须先从旧所有者的数据结构移除，再由新所有者解释字段。

\section{一次精确 buddy 演算：地址相邻不等于可以合并}

设页大小为 4096 字节，一个 zone 的 PFN 0--15 全部空闲。R37 的对象页最终只需 order 0，
但为展示拆分，先请求 order 1，即两个连续页。分配器从 order 4 的 PFN 0--15 块开始：

\begin{enumerate}
\tightlist
\item 拆成 PFN 0--7 与 8--15 两个 order 3 块，把 8--15 放回空闲表；
\item 把 0--7 拆成 0--3 与 4--7 两个 order 2 块，把 4--7 放回；
\item 把 0--3 拆成 0--1 与 2--3 两个 order 1 块，返回 0--1；
\item 页描述符把 PFN 0 标为已分配块头，PFN 2、4、8 分别记录空闲块阶数。
\end{enumerate}

\begin{pdfdiagram}[scale=0.9]
  \node[os state, minimum width=48mm] (o4) at (0,28mm)
    {order 4\\PFN 0--15};
  \node[os compute, minimum width=40mm] (o3l) at (-30mm,8mm)
    {order 3\\PFN 0--7};
  \node[os memory, minimum width=40mm] (o3r) at (30mm,8mm)
    {order 3 空闲\\PFN 8--15};
  \node[os compute, minimum width=34mm] (o2l) at (-48mm,-13mm)
    {order 2\\PFN 0--3};
  \node[os memory, minimum width=34mm] (o2r) at (-12mm,-13mm)
    {order 2 空闲\\PFN 4--7};
  \node[os control, minimum width=30mm] (alloc) at (-58mm,-34mm)
    {order 1 已分配\\PFN 0--1};
  \node[os memory, minimum width=30mm] (free) at (-25mm,-34mm)
    {order 1 空闲\\PFN 2--3};
  \draw[os bus] (o4.south) -- (o3l.north);
  \draw[os bus] (o4.south) -- (o3r.north);
  \draw[os bus] (o3l.south) -- (o2l.north);
  \draw[os bus] (o3l.south) -- (o2r.north);
  \draw[os strong bus] (o2l.south) -- (alloc.north);
  \draw[os bus] (o2l.south) -- (free.north);
\end{pdfdiagram}
\diagramnote{每次拆分只产生一个继续向下的候选块和一个放回相应 order 的空闲 buddy。释放 PFN 0--1 时，只有 PFN 2--3 仍为空闲 order 1 块才可以先合成 PFN 0--3。}

order \(k\) 块的 buddy 块头可以教学化地写成
\(\textit{buddyPFN}=\textit{blockPFN}\mathbin{\mathrm{xor}}2^k\)。
例如 PFN 0 的 order 1 buddy 是 PFN 2；PFN 4 虽与 PFN 2--3 相邻，却属于另一个 order 2
父块，不能与 PFN 0--1 直接合并。

\begin{lstlisting}
free_pages(block_pfn, order):
  // 只有块头记录当前 order；调用者必须使用与分配相同的阶数释放。
  while order < MAX_ORDER:
    buddy_pfn = block_pfn xor (1 << order)
    buddy = page_descriptor(buddy_pfn)

    // buddy 必须空闲、阶数相同、属于同一 zone 和迁移兼容集合。
    if not buddy.is_free or buddy.order != order:
      break
    if buddy.zone != page_descriptor(block_pfn).zone:
      break

    remove_from_free_list(buddy, order)
    block_pfn = min(block_pfn, buddy_pfn)
    order += 1

  // 最终合并块只登记一次，内部页不能同时出现在其他空闲块中。
  insert_free_block(block_pfn, order)
\end{lstlisting}

\subsection{空闲页守恒之外还要验证不重叠}

只比较分配前后的空闲页总数，无法发现同一页面被两个 order 同时登记。buddy 的核心不变量是：

\begin{itemize}
\tightlist
\item 每个空闲页恰好属于一个空闲块；
\item 每个空闲块按自身大小对齐，且不跨 zone 边界；
\item 已登记块之间不重叠；
\item 可合并的同阶 buddy 不应长期同时留在两个空闲表中；
\item 分配块只由块头携带阶数，释放接口必须取得正确的原始 order。
\end{itemize}

错误 order 的释放尤其危险。把一个 order 0 页面误报为 order 2，会让分配器把仍属于其他
对象的相邻页面登记为空闲；后续分配成功反而会产生两个合法调用者写同一物理页的破坏。

\section{zone、水位与 zonelist：有空闲页为何仍可能拒绝分配}

一个 zone 不能把最后几页全部交给普通可等待请求。中断、页表建立、回收本身和设备完成路径
仍可能需要内存；如果这些路径也因内存不足而无法前进，系统就不能释放任何页面。zone 因此
设置最低、较低和较高水位，并为原子上下文保留一部分应急容量。

\begin{longtable}{|L{0.18\linewidth}|L{0.3\linewidth}|L{0.42\linewidth}|}
\hline
状态 & 典型判定 & 分配与回收动作 \\ \hline
高于 high & 空闲量充足且目标 order 可取得 & 快路径分配，后台回收可停止 \\\tablerowrule
介于 low 与 high & 仍可服务普通请求，但余量下降 & 唤醒后台回收线程 \\\tablerowrule
介于 min 与 low & 普通分配需要谨慎推进 & 进入直接回收或按策略回退节点 \\\tablerowrule
低于 min & 普通请求不得继续消耗保留页 & 仅允许符合资格的原子/应急请求 \\\tablerowrule
总量足够但高阶不足 & 空闲页分散，watermark 可能通过 & 尝试 compaction 或降低 order \\
\hline
\end{longtable}

zonelist 表示一次请求允许尝试的 zone 和 NUMA 节点顺序。它不是固定的全机优先级：内存策略、
CPU 所在节点、地址限制和请求标志会共同决定候选顺序。分配器必须区分“本地节点没有合适块”
与“所有允许节点都失败”，否则既可能过早 OOM，也可能为一次小分配付出长期远端访问代价。

\begin{lstlisting}
alloc_from_zonelist(order, context):
  // 先按 NUMA 策略和地址限制生成本次请求可使用的 zone 顺序。
  candidates = build_zonelist(context.node_policy, context.zone_limit)

  for zone in candidates:
    // watermark 检查同时考虑 order、保留页和本请求是否有资格使用储备。
    if zone_watermark_ok(zone, order, context.reserve_class):
      block = remove_suitable_free_block(zone, order, context.migratetype)
      if block != NONE:
        return block

  // 快路径没有交付对象；是否回收由上下文决定，不能在这里无条件睡眠。
  return SLOWPATH_REQUIRED
\end{lstlisting}

\section{slab 几何：336 字节对象怎样落到页面中的槽位}

R37 的有效字段为 336 字节，要求 64 字节对齐。若调试配置在对象前后各放 16 字节红区，
并增加 16 字节元数据，单槽原始长度为 384 字节；向 64 字节对齐后仍是 384 字节。
一个 4096 字节页可排列 10 个槽，留下 256 字节给页尾空隙或其他 slab 元数据。调试关闭时，
实现可能选择不同布局；稳定语义是返回地址满足对象大小与对齐，而不是固定“每页十个”。

\begin{pdfdiagram}[scale=0.9]
  \node[os title] at (-61mm,21mm) {一个承载 R37 类型对象的 slab 页};
  \foreach \x/\lab/\sty in {
    -54/0/os state,
    -42/1/os state,
    -30/2/os memory,
    -18/3/os state,
    -6/4/os memory,
    6/5/os state,
    18/6/os state,
    30/7/os memory,
    42/8/os state,
    54/9/os memory} {
      \node[\sty, minimum width=10mm, minimum height=13mm, inner sep=1pt]
        at (\x mm,6mm) {\lab};
  }
  \node[os label, text width=45mm] at (-34mm,-12mm)
    {绿色：已分配对象\\黄色：freelist 中的空闲槽};
  \node[os label, text width=45mm] at (34mm,-12mm)
    {每槽 384 B\\页内余量 256 B};
\end{pdfdiagram}
\diagramnote{固定大小槽把任意字节切分转换为“从 freelist 取一个对象”。颜色表达槽位状态，不表达对象引用是否已经安全归零；发布和回收仍由对象所属子系统证明。}

\subsection{slab 页、对象槽和业务对象是三种身份}

\begin{longtable}{|L{0.2\linewidth}|L{0.3\linewidth}|L{0.4\linewidth}|}
\hline
身份 & 代表字段 & 所有权边界 \\ \hline
cache & name、object size、alignment、flags、constructor &
描述一种对象布局，跨越许多 slab 页 \\\tablerowrule
slab 页 & freelist、inuse、objects、node、list state &
由 cache 管理，在 partial/full/empty 之间迁移 \\\tablerowrule
对象槽 & slot address、debug state、freelist link &
空闲时由分配器拥有，分配后由调用子系统拥有 \\\tablerowrule
R37 业务对象 & id、generation、refcount、request state、completion &
由块层和驱动引用，不能用 slab 的 inuse 代替业务引用计数 \\
\hline
\end{longtable}

slab 的 \code{inuse} 只说明多少槽已经交给调用者，不知道 R37 是否仍在设备队列、等待队列
或完成路径中。业务对象必须先解除这些引用，才能调用释放接口。反过来，R37 的引用计数归零
也不负责把整页交回 buddy；cache 会根据 partial/empty 策略和压力决定何时释放 slab 页。

\Needspace{8\baselineskip}
\section{SLUB 快路径与慢路径：每 CPU freelist 怎样转移所有权}

多核同时访问一条全局 freelist 会让小对象分配受锁竞争限制。SLUB 类实现让当前 CPU 暂时
拥有一个活动 slab 和其中的 freelist。常见分配只弹出一个本地槽；本地耗尽以后，慢路径从
节点 partial 集合取得另一个 slab，或向 buddy 申请新页。

\begin{lstlisting}
allocate_r37(cache, context):
  // 1. 快路径只操作当前 CPU 已拥有的 freelist，不进入回收。
  local = cache.per_cpu[current_cpu()]
  obj = local.pop_with_consistency_check()
  if obj != NONE:
    initialize_r37(obj)
    return obj

  // 2. 慢路径先从本 NUMA 节点的 partial 集合转移一个 slab。
  slab = take_partial_slab(cache.node[context.preferred_node])
  if slab != NONE:
    local.install(slab)
    obj = local.pop_with_consistency_check()
    initialize_r37(obj)
    return obj

  // 3. 没有 partial slab 时才向页分配器请求承载页。
  slab = new_slab_from_pages(cache, context)
  if slab == NONE:
    return ALLOCATION_FAILED
  local.install(slab)
  obj = local.pop_with_consistency_check()
  initialize_r37(obj)
  return obj
\end{lstlisting}

“每 CPU”表示所有权和同步域，不表示对象永远由同一 CPU 释放。R37 可以在 CPU 2 创建，
由 CPU 7 的中断完成。远端释放要按实现协议安全地把槽接回目标 freelist 或节点 partial
集合，不能无锁修改另一个 CPU 正在消费的链表。

\subsection{CPU 下线时必须归还本地库存}

若 CPU 7 下线而本地 cache 仍持有 80 个空闲 R37 槽，这些槽仍计入系统内存，却不再有正常
分配路径访问。CPU 热插拔协议必须停止该 CPU 的分配活动、排空本地 freelist、把 partial
slab 归还节点集合，并在所有并发引用收束后释放 per-CPU 状态。

内存压力下也会主动 drain 每 CPU 缓存。这样做增加同步和缓存失效，却能把分散在各 CPU 的
空闲对象集中到 slab 页，使完全空闲的页有机会退回 buddy。快路径优化因此产生一项维护债务：
系统必须在 CPU 下线和回收时重新集中局部库存。

\section{分配慢路径：回收、压缩与 OOM 是不同失败的处理}

快路径失败后，系统先判断请求允许做什么。order 0 请求通常关心总量；高阶请求还关心连续性。
直接回收尝试释放可回收页面，compaction 尝试移动可移动页以形成连续区间，OOM 选择终止
某个资源消费者。三者处理的条件不同，不能把“再试一次”写成没有边界的循环。

\begin{pdfdiagram}[scale=0.88]
  \node[os state, minimum width=34mm] (request) at (-55mm,24mm)
    {请求\\order、zone、context};
  \node[os compute, minimum width=34mm] (fast) at (-18mm,24mm)
    {快路径\\watermark 与 free area};
  \node[os memory, minimum width=34mm] (reclaim) at (20mm,24mm)
    {直接回收\\扫描可回收页};
  \node[os control, minimum width=34mm] (compact) at (57mm,24mm)
    {连续性仍不足\\执行 compaction};
  \node[os state, minimum width=36mm] (success) at (20mm,-12mm)
    {重新检查成功\\返回分配块};
  \node[os control, minimum width=36mm] (oom) at (57mm,-12mm)
    {策略允许时\\进入 OOM 决策};
  \node[os memory, minimum width=36mm] (fail) at (-18mm,-12mm)
    {不可等待上下文\\立即失败或使用预留};
  \draw[os bus] (request.east) -- (fast.west);
  \draw[os bus] (fast.east) -- (reclaim.west);
  \draw[os bus] (reclaim.east) -- (compact.west);
  \draw[os strong bus] (reclaim.south) -- (success.north);
  \draw[os strong bus] (compact.south) -- (oom.north);
  \draw[os feedback] (fast.south) -- (fail.north);
\end{pdfdiagram}
\diagramnote{图中三条离开快路径的路线由分配上下文和失败类型选择：不可等待请求不能进入直接回收；总量改善后可直接成功；总量存在但连续性不足时才需要 compaction。}

\begin{longtable}{|L{0.2\linewidth}|L{0.27\linewidth}|L{0.31\linewidth}|L{0.12\linewidth}|}
\hline
阶段 & 输入证据 & 成功条件 & 可能代价 \\ \hline
后台回收 & zone 低于 low watermark & 空闲量回到目标水位 & 写回与缓存淘汰 \\\tablerowrule
直接回收 & 当前请求允许睡眠与回收 & 释放足够页面后重试成功 & 调用延迟不可预测 \\\tablerowrule
compaction & 高阶块不足且存在可移动页 & 形成目标 order 连续空闲块 & 迁移复制与 TLB 更新 \\\tablerowrule
memcg 回收 & 资源组 charge 达到限制 & 在同一 memcg 内释放容量 & 局部抖动 \\\tablerowrule
OOM 决策 & 允许回收路径均无法前进 & 终止对象释放足够资源 & 业务失败与恢复时间 \\
\hline
\end{longtable}

\subsection{compaction 的前后状态必须同时满足映射正确}

设 PFN 100--107 中，100、102、104、106 空闲，其他页可移动。系统总计有四个空闲页，却没有
order 2 连续块。compaction 为已占用页分配目标页，复制内容，更新所有映射与反向映射，
完成必要的 TLB 协议后才释放旧页。只复制字节而未更新引用会让访问继续指向旧 PFN。

\begin{lstlisting}
migrate_page(old_page, new_page):
  // 1. 隔离旧页，阻止建立新的不受控引用；不可移动或被 pin 时立即失败。
  if not isolate_movable_page(old_page):
    return NOT_MOVABLE

  // 2. 锁定并复制内容；复制期间旧映射不能观察半更新状态。
  lock(old_page)
  copy_page_bytes(old_page, new_page)
  copy_required_page_state(old_page, new_page)

  // 3. 把文件/匿名映射、页表项和反向映射切换到新页。
  replace_all_mappings(old_page, new_page)
  complete_required_tlb_invalidation()

  // 4. 新映射已经生效后，旧页才重新进入空闲集合。
  unlock(old_page)
  free_migrated_source(old_page)
  return MIGRATED
\end{lstlisting}

被设备长期 pin 的页、部分内核页、不可迁移页表状态或处于写回冲突的页会阻断压缩。诊断
高阶失败时必须同时观察不可移动页分布和迁移失败原因；只报告“尚有 20\% 空闲内存”不能
解释连续块为何不存在。

\section{分配上下文：GFP 规则是在阻止递归等待}

R37 的普通提交路径可以使用允许等待的上下文；设备中断完成路径不能睡眠。更隐蔽的问题是
“可以睡眠”不等于“可以进入任意回收”。若文件系统持有 inode 锁时分配内存，直接回收又
进入同一文件系统写回并等待 inode 锁，原任务会等待自己释放锁。

\begin{longtable}{|L{0.2\linewidth}|L{0.26\linewidth}|L{0.32\linewidth}|L{0.12\linewidth}|}
\hline
上下文 & 允许动作 & 禁止或受限原因 & 失败策略 \\ \hline
普通进程路径 & 睡眠、直接回收、按契约启动 I/O & 仍需服从已持锁和递归边界 & 返回错误或 OOM \\\tablerowrule
文件系统递归敏感路径 & 睡眠、非文件系统回收 & 重新进入文件系统可能等待当前锁 & 预留、延迟或失败 \\\tablerowrule
I/O 递归敏感路径 & 睡眠、避免发起新 I/O 的回收 & 新 I/O 可能依赖当前请求完成 & 预留或失败 \\\tablerowrule
中断/自旋锁路径 & 非阻塞快路径与授权预留 & 调度睡眠会破坏执行上下文 & 丢弃、延迟工作或预分配 \\\tablerowrule
恢复关键路径 & 预留池中保证的最小对象 & 失败会阻止释放资源本身 & 消费 mempool，完成后归还 \\
\hline
\end{longtable}

\subsection{mempool 解决的是前进保证，不是无限容量}

若块层错误恢复至少需要 16 个请求对象才能排空已提交队列，系统可以预先建立 16 个 R37
预留对象。正常分配先用普通 cache；普通路径失败时，符合资格的恢复路径消费预留。对象完成
后优先补回预留池，恢复其最小保证。

\begin{lstlisting}
alloc_recovery_request(pool, context):
  // 正常路径不消耗稀缺预留，先尝试普通对象 cache。
  obj = slab_try_alloc(pool.cache, context.without_recursive_reclaim)
  if obj != NONE:
    return obj

  // 只有声明为恢复关键的调用者可以消费预留对象。
  if context.recovery_critical:
    obj = pool.take_reserved()
    if obj != NONE:
      return obj

  // 预留也耗尽时必须显式失败，不能在不可等待路径循环。
  return ALLOCATION_FAILED
\end{lstlisting}

预留池的大小来自可证明的最小并发需求。若所有普通调用都能消费它，真正的恢复路径仍会在
压力峰值时失败；若预留过大，则大量内存长期不可供其他请求使用。

\section{接口选择：连续的究竟是地址、PFN 还是设备视图}

“需要一块连续缓冲区”必须说明由谁观察连续。CPU 可以通过页表把离散物理页映射成连续虚拟
区间；支持 scatter-gather 的设备可以通过描述符看到多个片段；IOMMU 还可以为设备建立连续
I/O 虚拟地址。只有明确观察者，才能选择接口。

\Needspace{18\baselineskip}
\begin{longtable}{|L{0.19\linewidth}|L{0.23\linewidth}|L{0.28\linewidth}|L{0.2\linewidth}|}
\hline
需求 & 合适来源 & 获得的连续性 & 不能据此推断 \\ \hline
\endfirsthead
\hline
需求 & 合适来源 & 获得的连续性 & 不能据此推断 \\ \hline
\endhead
固定小对象 R37 & 专用 slab cache & 单对象虚拟地址与所需对齐 & 多对象彼此物理相邻 \\\tablerowrule
若干完整连续页 & buddy 高阶分配 & PFN 物理连续且内核可映射 & 高阶请求必定成功 \\\tablerowrule
大型 CPU 缓冲 & vmalloc 类映射 & 内核虚拟地址连续 & 可直接用于无 SG 的 DMA \\\tablerowrule
设备 DMA 区域 & DMA API、SG 或一致性分配 & 符合设备和 IOMMU 契约的设备视图 & CPU 虚拟地址等于 DMA 地址 \\\tablerowrule
紧急对象集合 & mempool 加底层 allocator & 规定数量的前进保证 & 无限并发时仍不失败 \\
\hline
\end{longtable}

R37 本体只是控制对象，使用专用 cache 最合适；它引用的数据页可以由页缓存或 DMA 映射单独
管理。把控制对象和大数据缓冲强制放进同一连续分配，会放大高阶失败并让两个不同生命周期
互相约束。

\section{对象发布与释放：refcount 归零以后仍可能不能复用}

R37 创建后经历 INIT、PUBLISHED、IN\_FLIGHT、COMPLETING 和 RETIRED。设备队列、等待者、
超时工作和查找表可能同时持有引用。把对象从查找表移除，只阻止未来查找；已经取得指针的
RCU 读者仍可能继续访问。因此“不可再找到”“引用计数归零”“宽限期结束”和“槽位回到
freelist”是四个边界。

\begin{longtable}{|L{0.18\linewidth}|L{0.28\linewidth}|L{0.34\linewidth}|L{0.1\linewidth}|}
\hline
阶段 & 所有者与代表字段 & 进入下一阶段的证据 & 可复用 \\ \hline
INIT & 创建线程；字段尚未公开 & 初始化完成，发布写建立顺序 & 否 \\\tablerowrule
PUBLISHED & 全局表和调用者持有引用 & 队列接收并取得自身引用 & 否 \\\tablerowrule
IN\_FLIGHT & 设备/驱动、等待者、超时工作 & 完成或复位证明设备不再访问 & 否 \\\tablerowrule
RETIRED & 已从查找表移除，旧读者可能存在 & refcount 归零并经过所需宽限期 & 否 \\\tablerowrule
QUARANTINED & 调试器持有，内容已 poison & 隔离策略允许重新进入 freelist & 否 \\\tablerowrule
FREE & slab freelist 拥有槽位 & 新分配原子取得所有权 & 是 \\
\hline
\end{longtable}

\begin{lstlisting}
retire_request(req):
  // 1. 先阻止新查找；已持有引用的读者仍可访问。
  remove_from_request_index(req.id)
  store_release(req.state, RETIRED)

  // 2. 释放索引持有的引用；最后引用可能在另一 CPU 的完成路径归零。
  if refcount_put(req.refcount) == LAST_REFERENCE:
    call_after_rcu_grace_period(final_free_request, req)

final_free_request(req):
  // 3. 宽限期证明旧 RCU 读者结束，再检查设备和调试状态。
  require(req.device_mapping == NONE)
  require(req.timer_state == CANCELLED)
  poison_and_quarantine_or_free(req)
\end{lstlisting}

\subsection{generation 防止旧完成命中新对象}

槽位释放后可能很快重新用于 R52，地址恰好与旧 R37 相同。若设备复位前的迟到完成只携带
槽位地址或短标签，它可能错误地把 R52 标为完成。请求身份要组合 tag 与 queue generation；
队列复位推进 generation，旧完成即使到达也只能结算为过期证据，不能修改新对象。

这项规则说明分配器能够复用地址，却不能替子系统复用身份。地址相同只是存储位置再次分配，
不表示旧异步事件属于新生命周期。

\section{资源归属与失败范围：memcg OOM 不等于整机 OOM}

R37 的控制对象、引用的数据页和间接分配可能属于不同计费类别。资源组 charge 要在对象对
业务可见前建立；若 charge 失败，初始化路径必须按相反顺序释放已经取得的页、槽和引用。
释放时 uncharge 只能发生一次，通常与最终资源所有权解除绑定，而不是与某个用户句柄关闭
简单等同。

\begin{lstlisting}
create_charged_request(owner, bytes):
  // 1. 先预留资源组额度；失败时还没有可见对象。
  charge = memcg_try_charge(owner.memcg, bytes)
  if charge == FAILED:
    return MEMCG_LIMIT

  // 2. 分配和初始化失败时必须撤销已经成功的 charge。
  req = allocate_r37(request_cache, owner.allocation_context)
  if req == ALLOCATION_FAILED:
    memcg_uncharge(charge)
    return NO_MEMORY

  // 3. 对象保存 charge 身份，最终释放路径负责恰好撤销一次。
  req.memcg_charge = charge
  initialize_owner_fields(req, owner)
  publish_request(req)
  return req
\end{lstlisting}

memcg 达到上限时，回收应优先限制在该资源组；无法前进时触发局部 OOM，不应因为一个容器
超限就任意终止宿主其他业务。只有全局允许 zone 都无法满足请求，问题才属于整机容量或碎片。

\section{调试证据：错误访问发生点与错误释放点往往不同}

对象越界可能在 R37 正常运行时破坏相邻槽，直到邻居释放才被发现；释放后访问可能在槽位
重新分配为另一个类型后才表现为随机字段损坏。调试机制需要分别覆盖边界、时间和可达性。

\Needspace{20\baselineskip}
\begin{longtable}{|L{0.18\linewidth}|L{0.25\linewidth}|L{0.31\linewidth}|L{0.16\linewidth}|}
\hline
证据机制 & 保存状态 & 能定位的问题 & 主要代价 \\ \hline
\endfirsthead
\hline
证据机制 & 保存状态 & 能定位的问题 & 主要代价 \\ \hline
\endhead
redzone/\allowbreak canary & 对象边界前后的固定模式 & 相邻越界写与部分下溢 & 增加槽大小 \\\tablerowrule
poison & 分配/释放时填入特征字节 & 未初始化读取、释放后内容变化 & 写入开销 \\\tablerowrule
quarantine & 延迟复用及释放栈 & 提高 UAF 在旧生命周期内暴露概率 & 内存滞留 \\\tablerowrule
shadow memory & 每段地址的可访问状态 & 精确越界与 UAF 访问点 & 较高内存/性能开销 \\\tablerowrule
guard page sampling & 少量对象邻接不可访问页 & 低开销捕获部分堆越界 & 非确定覆盖 \\\tablerowrule
alloc/free stack & 调用点、CPU、时间和 generation & 把错误访问连接到生命周期 & 记录与存储成本 \\
\hline
\end{longtable}

检测到 canary 损坏时，报告应包含 cache 名称、对象地址、slab 页、相邻槽状态、分配栈、
释放栈、CPU、generation 和首次异常字节偏移。只打印“内存损坏”不能区分分配器元数据错误、
调用者越界还是错误 order 释放。

\section{完整复盘：R37 怎样从物理页中取得身份并安全消失}

现在把开篇任务逐层闭合：

\begin{enumerate}
\item 启动阶段从固件候选 RAM 中扣除内核、页表和保留范围，为 managed PFN 建立页描述符，
      再按 zone、order 和迁移类型放入 buddy。
\item 请求 cache 保存 R37 的有效大小、64 字节对齐、构造与调试布局。首次缺少 slab 页时，
      cache 通过允许等待的上下文向 buddy 取得页面并切成固定槽。
\item 当前 CPU 从本地 freelist 取得一个槽；慢路径可能转移 partial slab 或新建 slab。
      槽位所有权从分配器转给块层，但此时 R37 尚未发布。
\item 创建线程初始化 id、queue generation、refcount、状态、等待队列和资源 charge；发布写
      完成后，全局索引和设备队列才可以取得引用。
\item 普通提交路径可以按上下文进入回收；中断或恢复关键路径只能使用非阻塞快路径和授权
      预留。失败结果必须返回调用者或执行明确降级，不能在不可等待路径隐式睡眠。
\item 完成路径先撤销设备访问能力，再发布结果并移除索引。refcount 归零只结束显式引用；
      RCU 宽限期和调试 quarantine 结束后，槽位才回到 freelist。
\item 当 slab 页全部为空且缓存策略允许回收时，页才从对象分配器退回 buddy。buddy 按相同
      order 验证并合并伙伴，最终恢复页级连续空闲块。
\end{enumerate}

这条路径把三组相邻概念分开：虚拟地址连续、物理页连续和设备地址连续；空闲页总量、最大
连续块和可用应急储备；对象不可查找、引用归零和内存可复用。读者能够指出每次所有权转移的
字段和证据以后，buddy、SLUB、GFP、mempool、compaction 与 memcg 才不再是孤立名称。

\Needspace{32\baselineskip}
\subsection{一次失败回放：为什么“还有空闲内存”不足以定位}

假设 R37 的 cache 补充 slab 时请求 order 2 页面失败，而监控仍显示 18\% 空闲内存。有效
诊断不能从空闲百分比直接跳到 OOM，至少要按下面的证据顺序回放：

\begin{enumerate}
\tightlist
\item 先核对请求的 order、允许 zone、NUMA 策略和 allocation context，确认它究竟在什么
      地址范围内寻找哪一种连续块，以及是否允许睡眠和回收；
\item 再读取候选 zone 的 watermark、各 order 空闲块计数和不可移动页分布，区分“总量低”
      与“总量分散”；
\item 若直接回收已经释放页面，记录释放页进入了哪个 zone 和迁移类型；释放到不合资格的
      区域不会使当前请求成功；
\item 若 compaction 失败，记录首个不可迁移 PFN、pin 来源、写回状态和迁移错误，而不是只
      保存一个笼统失败码；
\item 最后才依据上下文决定返回失败、使用预留、降低 order 或进入 OOM，并把所选分支写入
      请求级追踪记录。
\end{enumerate}

这份记录把“调用者提出了什么契约”“分配器看到了什么状态”和“策略为何选择该结果”连在
一起。只有三类证据能在同一时间线上对齐，才可能判断应当调整对象布局、减少长期 pin、改变
NUMA 策略，还是修复错误的分配上下文。
